\documentclass[UTF8,12pt,a4paper,fontset=fandol]{ctexart}

% 支持从项目根目录、深度学习目录或当前专题目录编译。
\makeatletter
\def\input@path{{../}{../../}}
\makeatother
\usepackage{researchnotes}

\title{低秩适配的原理与数学基础}
\author{}
\date{}

\begin{document}
\maketitle

\begin{abstract}
低秩适配（low-rank adaptation, LoRA）通过冻结预训练权重、训练低秩增量，降低任务适配所需的可训练参数量。理解这一方法，需要同时掌握矩阵的秩与分解、监督学习的目标函数、反向传播、优化器状态以及大模型的计算结构。本文从可手算的线性变换出发，建立这些前置知识，推导低秩分支的前向计算、梯度、初始化和训练动力学，分析参数效率、表达限制、显存构成、量化与权重合并，并给出实验设计、故障诊断和带解答的练习。全文区分数学恒等式、附加假设下的分析和需要实验验证的经验判断；不把小秩有效解释为普适定理，也不把参数压缩比例当作训练加速比例。
\end{abstract}

\tableofcontents

\section{从一个具体问题理解微调与低秩适配}
\label{sec:motivation}

\subsection{预训练模型已经会什么，任务适配还要改变什么}

设一个模型已通过大量数据学会一般语言规律，现在希望它把客服问题分为退款、物流和账户三个类别。一种方案是更新模型中的全部参数，另一种方案是保留原有参数，只学习能够改变任务行为的少量附加参数。前者称为\keyterm{全参数微调}（full fine-tuning），后者属于\keyterm{参数高效微调}（parameter-efficient fine-tuning, PEFT）。这里的“高效”首先指需要训练和保存的任务参数较少，并不预先承诺每个任务的准确率更高，也不承诺所有硬件上都有同样的速度收益。

记预训练参数为 $\theta_0$，微调后的参数为 $\theta_0+\Delta\theta$。全参数微调通常允许 $\Delta\theta$ 在全部参数坐标上变化，低秩适配则对其中某些权重矩阵的增量施加结构限制。它的出发点是：已有模型可能提供了足够丰富的基础表示，完成一个新任务不一定需要任意改变所有权重。这个判断是方法设计的动机；某个任务究竟需要多少自由度，应通过数据、基线和验证实验判断，不能由“模型很大”直接推出“任意微调都是低秩的”。

\subsection{先看一个秩一修正}

采用列向量约定，线性层将 $x\in\mathbb R^3$ 映射为 $h=W_0x\in\mathbb R^2$。令
\[
W_0=\begin{pmatrix}1&0&0\\0&1&0\end{pmatrix},
\qquad
A=\begin{pmatrix}1&-1&0\end{pmatrix},
\qquad
B=\begin{pmatrix}2\\1\end{pmatrix}.
\]
加入修正后得到 $h=W_0x+BAx$。矩阵 $A$ 先从输入中提取一个标量 $x_1-x_2$，矩阵 $B$ 再决定如何把这个标量分配给两个输出：第一个输出增加其两倍，第二个输出增加其一倍。因此
\[
BA=\begin{pmatrix}2&-2&0\\1&-1&0\end{pmatrix},
\qquad
(W_0+BA)\begin{pmatrix}3\\1\\4\end{pmatrix}
=\begin{pmatrix}7\\3\end{pmatrix}.
\]
这不是删去原模型的第三个输入，也不是把整层变成只有一个输出的网络；原路径 $W_0x$ 仍然存在，新增路径只根据一个可学习的组合方向实施修正。第三个输入在本例中不影响输出，是所选矩阵的特殊性质，不是低秩适配的一般要求。

上述修正有六个矩阵元素，却可用 $A$ 的三个数和 $B$ 的两个数表示。这个小例子的节省很有限，大矩阵才更能体现收益。若权重为 $4096\times4096$，直接训练它需要 $16777216$ 个标量；用中间维度 $r=8$ 的两个因子表示增量，只需 $8(4096+4096)=65536$ 个标量。与此同时，原来的大矩阵仍需加载并参与计算。理解“训练对象变小”和“整个模型变小”的区别，是阅读后文的第一条主线。

\subsection{全文的知识依赖与研究边界}

后文首先解释向量空间、矩阵秩、奇异值分解和低秩逼近，再从概率模型建立训练损失，用矩阵微分推导梯度，并说明 Transformer 中适合插入修正的位置。在这些基础上，讨论 LoRA 的参数化、初始化、优化几何和秩选择，然后分析显存、量化、实现与评估。最后讨论合并、组合和代表性变体，并用习题检查理解。LoRA 的基本方案见原始论文\cite{hu2021}；本文的算例和推导用于教学，不对应未经说明的模型实测结果。

\section{线性代数前置知识：矩阵、子空间与秩}
\label{sec:linear-algebra}

\subsection{矩阵是一个线性映射}

向量空间 $\mathbb R^n$ 中的向量 $x=(x_1,\ldots,x_n)^{\mathsf T}$ 可表示一个词元的隐藏特征。这里 $n$ 是特征维度，不是词元数量。矩阵 $W\in\mathbb R^{m\times n}$ 定义线性映射 $x\mapsto Wx$，其中第 $i$ 个输出为 $\sum_{j=1}^{n}W_{ij}x_j$。若增加偏置 $b\in\mathbb R^m$，得到的是仿射映射 $Wx+b$。数学讨论中常暂时省略偏置，实际实现却必须明确它是冻结、训练还是单独保存的参数。

将 $W$ 按列写为 $[w_1,\ldots,w_n]$，则 $Wx=\sum_jx_jw_j$。因此所有可能的输出都位于这些列向量张成的空间中。一个矩阵有很多非零元素，不代表这些元素描述了同样多的独立作用方向；例如全部列都互为倍数时，输入无论如何变化，输出都只能沿一条直线移动。

\begin{definition}[秩、像空间与零空间]
对 $W\in\mathbb R^{m\times n}$，其像空间为
$\operatorname{im}(W)=\{Wx:x\in\mathbb R^n\}$，其维数称为矩阵的秩（rank），记为 $\operatorname{rank}(W)$。其零空间为
$\ker(W)=\{x:Wx=0\}$。秩等于列空间的维数，也等于行空间的维数，并满足
$0\leq\operatorname{rank}(W)\leq\min(m,n)$。
\end{definition}

秩与参数个数不同，秩与稀疏性也不同。\keyterm{稀疏矩阵}（sparse matrix）含有大量零元素，而低秩矩阵可以每个元素都非零。例如全一矩阵的秩为一；单位矩阵只有对角线上非零，却是满秩。LoRA 的“低秩”限制独立变化方向的数量，不要求增量中只有少量位置发生变化，因此它能够用较少参数同时改变大量权重元素。

\subsection{矩阵乘积为什么产生低秩结构}

取 $A\in\mathbb R^{r\times n}$、$B\in\mathbb R^{m\times r}$。先计算 $z=Ax$，再计算 $Bz$，相当于让新增变化经过一个 $r$ 维中间表示。无论输入怎样变化，$BAx$ 都属于 $B$ 的列空间，因此有
\begin{equation}
\operatorname{rank}(BA)\leq
\min\{\operatorname{rank}(A),\operatorname{rank}(B)\}\leq r.
\label{eq:rank-product}
\end{equation}
这给出了低秩因子化的基本保证。反过来，任何秩不超过 $r$ 的 $m\times n$ 矩阵都能写成这样的乘积：选择其列空间的一组基放入 $B$，把各列关于这组基的坐标放入 $A$，不足 $r$ 的部分补零即可。

进一步把 $B$ 的第 $k$ 列记为 $b_k$，把 $A$ 的第 $k$ 行记为 $a_k^{\mathsf T}$，则
\begin{equation}
BA=\sum_{k=1}^{r}b_ka_k^{\mathsf T},
\qquad
BAx=\sum_{k=1}^{r}b_k(a_k^{\mathsf T}x).
\label{eq:outer-product}
\end{equation}
每一项都是“检测输入方向，再沿输出方向施加变化”。这些方向通常通过训练共同形成，不能先验地把某个方向命名为语法、情感或知识。中间坐标可以重新组合而不改变乘积，因此对单个因子元素的语义解释需要格外谨慎。

\subsection{实际秩、设定秩与可表示自由度}

在工程配置中，“LoRA 的秩为 $r$”通常是说中间维度设为 $r$，即增量秩的上界为 $r$。如果 $A$ 的行相关、$B$ 的列相关，或者某些方向的权重接近零，乘积的实际秩可能更低。浮点计算还需要区分精确秩与\keyterm{数值秩}（numerical rank）：后者依赖一个阈值，把足够小的奇异值视为零，所以报告数值秩时必须同时报告判定标准。

两个因子共有 $r(m+n)$ 个参数，但这些参数并非相互独立的有效自由度。对任意可逆矩阵 $C\in\mathbb R^{r\times r}$，
\[
(BC)(C^{-1}A)=BA.
\]
在恰好秩为 $r$ 的非退化情形，秩 $r$ 矩阵集合的局部维数为 $r(m+n-r)$。直观上，可以从两个因子的参数量中扣除中间基变换的 $r^2$ 个冗余方向；严格的局部坐标论证需要流形知识，这里只用这一计数说明“存储了多少参数”与“表示了多少不同矩阵”不是同一问题。对低于 $r$ 的秩层，该集合还包含不同维度的部分，不能把它整体当作同一光滑空间。

\subsection{范数、内积与尺度}

对同形矩阵 $U,V$，Frobenius 内积定义为
$\langle U,V\rangle_F=\operatorname{tr}(U^{\mathsf T}V)=\sum_{ij}U_{ij}V_{ij}$，其中 $\operatorname{tr}$ 是对角元素之和。Frobenius 范数为 $\|U\|_F=\sqrt{\langle U,U\rangle_F}$，可衡量全部元素的总体大小。谱范数定义为
$\|U\|_2=\max_{\|x\|_2=1}\|Ux\|_2$，衡量最容易被放大的输入方向。二者都能描述更新幅度，但回答的问题不同。

小秩并不蕴含小范数。矩阵 $10^6uv^{\mathsf T}$ 即使秩为一，也可能造成很大的输出变化。反之，一个满秩矩阵乘上极小系数后，其作用可能很弱。后续讨论中的秩 $r$、缩放系数和学习率分别控制结构、前向幅度与优化步长，不能把三个概念混成一个“适配强度”。

\section{奇异值分解与低秩逼近的适用边界}
\label{sec:svd}

\subsection{奇异值分解把矩阵拆成独立方向}

对实矩阵 $M\in\mathbb R^{m\times n}$，\keyterm{奇异值分解}（singular value decomposition, SVD）给出
\[
M=U\Sigma V^{\mathsf T}
=\sum_{j=1}^{q}\sigma_ju_jv_j^{\mathsf T},
\qquad q=\min(m,n),
\]
其中 $U,V$ 是相应的正交矩阵，$\sigma_1\geq\cdots\geq\sigma_q\geq0$，$u_j,v_j$ 分别是左、右奇异向量。几何上，右奇异向量给出输入的正交方向，奇异值给出伸缩强度，左奇异向量给出对应输出方向。这里引用实矩阵 SVD 的存在性，不在本文重证谱定理。

奇异值可以由半正定矩阵 $M^{\mathsf T}M$ 的特征值开平方得到，因此均为非负数。非零奇异值的个数等于秩，并且
$\|M\|_F^2=\sum_j\sigma_j^2$、$\|M\|_2=\sigma_1$。若只有前几个奇异值较大，矩阵虽然在精确意义上满秩，却可能具有良好的低秩近似。这种谱集中为低秩表示提供直觉，但应说明讨论的是哪一个矩阵：原始权重、梯度、训练后的增量，还是任务数据上的输出变化。

\begin{theorem}[截断奇异值分解的最小平方误差]
\label{thm:truncated-svd}
设 $1\leq r<q$，令 $M_r=\sum_{j=1}^{r}\sigma_ju_jv_j^{\mathsf T}$。则
\[
\min_{\operatorname{rank}(D)\leq r}\|M-D\|_F^2
=\|M-M_r\|_F^2=\sum_{j=r+1}^{q}\sigma_j^2.
\]
\end{theorem}

\begin{proof}
取任意秩不超过 $r$ 的 $D$，令 $P$ 为投影到其列空间的正交投影矩阵。由 $PD=D$ 以及正交分解，
$\|M-D\|_F^2\geq\|(I-P)M\|_F^2$。将 $MM^{\mathsf T}$ 用左奇异向量对角化，可得
$\|PM\|_F^2=\operatorname{tr}(PMM^{\mathsf T})
=\sum_j\sigma_j^2\|Pu_j\|_2^2$。各系数位于 $[0,1]$，总和不超过 $r$，所以该加权和至多为 $\sum_{j=1}^r\sigma_j^2$：在系数总量固定时，把系数优先分配给较大奇异值只会增大和。于是
$\|(I-P)M\|_F^2\geq\sum_{j>r}\sigma_j^2$。取 $D=M_r$ 达到下界，结论成立。
\end{proof}

例如 $M=\operatorname{diag}(5,2,1)$ 的最佳秩一近似可取
$\operatorname{diag}(5,0,0)$，平方误差为 $2^2+1^2=5$；最佳秩二近似的平方误差为 $1$。若临界奇异值存在重数，最优方向未必唯一。误差由被截去的谱决定，而不是由矩阵大小单独决定，所以“取固定小秩总能逼近大矩阵”没有数学依据。

\subsection{为什么 LoRA 不等于先训练再做奇异值分解}

截断 SVD 解决的是：给定一个已知矩阵 $M$，寻找 Frobenius 范数下最接近它的低秩矩阵。LoRA 解决的则是：在网络损失函数下直接训练两个因子，使任务表现改善。训练前并不知道全参数微调会得到哪个增量，而且不同学习率、随机种子和正则化可能产生不同增量。即便先得到某个全参数解，最接近其权重的低秩矩阵也未必是任务损失最小的低秩矩阵。

这个区别可以直接通过输入分布说明。设 $D$ 是拟合误差矩阵，随机输入 $X$ 满足有限二阶矩，记 $C_X=\mathbb E[XX^{\mathsf T}]$。注意 $C_X$ 是二阶矩矩阵，只有 $\mathbb E[X]=0$ 时才等于协方差矩阵。于是
\begin{equation}
\mathbb E\|DX\|_2^2
=\operatorname{tr}(DC_XD^{\mathsf T})
=\|DC_X^{1/2}\|_F^2.
\label{eq:data-weighted-error}
\end{equation}
第一个等式可逐元素展开后用有限求和与期望的线性性证明；有限二阶矩保证各项可积。只有当 $C_X$ 与单位矩阵成比例时，该误差才与普通 Frobenius 误差成比例。若某个输入方向在任务数据中几乎不出现，沿该方向产生较大权重误差也可能影响很小。

\subsection{一个数据加权误差的反例}

令目标增量为 $M=\operatorname{diag}(3,2)$，输入二阶矩为
$C_X=\operatorname{diag}(1,100)$。按普通 SVD 保留第一维，得到 $D_1=\operatorname{diag}(3,0)$，其任务输出平方误差为 $2^2\cdot100=400$；保留第二维，得到 $D_2=\operatorname{diag}(0,2)$，误差却只有 $3^2\cdot1=9$。二者都是秩一，但权重上最接近目标的方案不是此数据分布下输出误差最小的方案。

在深层非线性网络中，后续层还会放大或抑制不同输出方向，所以仅检查某层增量的奇异值不足以确定任务效果。SVD 是理解与诊断的工具，不能替代训练目标。实践中可以报告前 $r$ 个奇异值的能量占比
$\sum_{j\leq r}\sigma_j^2/\sum_j\sigma_j^2$，但需在分母非零时定义，并说明它描述矩阵能量而非准确率占比。

\section{概率、损失函数与微调究竟优化什么}
\label{sec:objective}

\subsection{总体目标与样本目标}

设 $(X,Y)$ 表示来自目标任务分布 $P$ 的随机输入和标签，$f_\theta$ 是参数为 $\theta$ 的模型，$\ell$ 是单样本损失。理想的总体风险为
$R(\theta)=\mathbb E_{(X,Y)\sim P}[\ell(f_\theta(X),Y)]$，但通常只能观察训练数据
$\mathcal D=\{(x_i,y_i)\}_{i=1}^N$，于是最小化经验风险
\begin{equation}
\widehat R(\theta)=\frac1N\sum_{i=1}^{N}
\ell(f_\theta(x_i),y_i).
\label{eq:empirical-risk}
\end{equation}
若样本独立同分布且损失可积，固定参数下样本平均可作为总体风险的估计。训练过程中参数依赖训练集，训练误差小并不自动意味着总体误差小，这正是必须使用独立验证集和测试集的原因。

LoRA 不改变这一区分。它把可优化参数从 $\theta$ 改为一组因子 $\phi$，并通过映射 $\theta(\phi)$ 得到实际网络权重；最终目标仍是 $\widehat R(\theta(\phi))$。任务可以是分类、序列生成或其他可微损失，并不因为使用 LoRA 就自动成为一种新的学习任务。监督微调（supervised fine-tuning, SFT）描述训练数据和监督方式，LoRA 描述参数化方式，二者可同时使用。

\subsection{从最大似然到交叉熵}

分类模型为每个类别给出概率 $p_\theta(k\mid x)$。对观测标签 $y$，负对数似然（negative log-likelihood, NLL）损失为
$-\log p_\theta(y\mid x)$。若模型先输出实数分数 $z_1,\ldots,z_K$，则软最大函数（softmax）定义
$p_k=e^{z_k}/\sum_je^{z_j}$。把真实标签写成独热向量 $q$，交叉熵（cross-entropy）为
$-\sum_kq_k\log p_k$，在独热标签下与负对数似然相同。标签平滑等方法会改变 $q$，因而也改变目标。

对固定真实条件分布 $q(\cdot\mid x)$，交叉熵可以分解为熵加 Kullback--Leibler 散度：
\[
-\sum_kq_k\log p_k
=-\sum_kq_k\log q_k+\sum_kq_k\log\frac{q_k}{p_k}.
\]
约定 $0\log0=0$，并要求 $q_k>0$ 的位置有 $p_k>0$，否则交叉熵和散度可能为无穷大。第一项不依赖模型，最小化交叉熵等价于在表达能力允许的范围内逼近真实分布；但有限数据、模型限制和优化误差使这个理想目标未必完全实现。

softmax 交叉熵对分数的导数为 $\partial\ell/\partial z_k=p_k-q_k$。例如两类概率为 $(0.7,0.3)$，真实类别为第二类，则导数为 $(0.7,-0.7)$，梯度下降倾向于降低第一类分数、提高第二类分数。梯度随后通过输出层和隐藏层传递到适配器，决定应改变哪些输入方向与输出方向。

\subsection{语言模型、词元与损失掩码}

自回归语言模型（autoregressive language model）将回答序列 $y=(y_1,\ldots,y_T)$ 的条件概率分解为
$p_\theta(y\mid x)=\prod_{t=1}^{T}p_\theta(y_t\mid x,y_{<t})$。其中词元（token）是分词器定义的离散单位，不必等于一个汉字、一个英文单词或一个字节。训练时提供真实前缀 $y_{<t}$ 的做法称为教师强制（teacher forcing）；生成时则使用模型此前生成的词元，两种输入分布可能不同。

对于批内样本，设 $m_{it}\in\{0,1\}$ 表示第 $i$ 条样本的第 $t$ 个位置是否纳入监督，且有效位置总数 $M=\sum_{i,t}m_{it}>0$，则可采用
\begin{equation}
\mathcal L=-\frac1M\sum_{i,t}m_{it}
\log p_\theta(y_{it}\mid x_i,y_{i,<t}).
\label{eq:token-loss}
\end{equation}
填充位置通常不应贡献损失；指令微调还需明确是否只监督回答部分。损失掩码决定哪些预测承担训练目标，注意力掩码决定哪些位置能看到哪些上下文，两者不能互相替代。

按有效词元平均与先按样本平均再对样本求平均通常不同：前者使长回答拥有更大权重，后者使每条样本拥有相同总权重。比较全参数微调与 LoRA 时，若两者采用不同的归一化、截断策略或监督范围，就无法把结果差异干净地归因于参数化方式。数据处理属于目标函数的一部分，不只是训练前的机械准备。

\section{反向传播、矩阵微分与优化器}
\label{sec:backprop}

\subsection{梯度是局部变化率}

对可微标量函数 $\mathcal L(W)$，矩阵梯度 $G=\nabla_W\mathcal L$ 由
$\mathrm d\mathcal L=\langle G,\mathrm dW\rangle_F$ 定义。若只改变一个元素，$G_{ij}$ 就是对应偏导数。对充分小的扰动 $D$，
$\mathcal L(W+D)=\mathcal L(W)+\langle G,D\rangle_F+o(\|D\|_F)$。
这是一阶局部近似，不意味着沿负梯度移动任意距离都能降低损失。

对 $h=Wx$，设上游梯度为 $g=\nabla_h\mathcal L\in\mathbb R^m$。由于
$\mathrm dh=(\mathrm dW)x+W\,\mathrm dx$，
\[
\mathrm d\mathcal L
=g^{\mathsf T}(\mathrm dW)x+g^{\mathsf T}W\,\mathrm dx,
\qquad
\nabla_W\mathcal L=gx^{\mathsf T},
\qquad
\nabla_x\mathcal L=W^{\mathsf T}g.
\]
外积 $gx^{\mathsf T}$ 描述单样本对线性权重的梯度，其秩至多为一。然而多个样本、多个词元和多个训练步骤的梯度相加后，秩可以增加到很大。因此“单样本梯度是秩一”只能提供局部结构直觉，不能证明整个微调增量一定低秩。

\subsection{链式法则与冻结参数}

反向传播（backpropagation）按照计算图反向应用链式法则，把损失对输出的导数逐层传回各参数。自动微分（automatic differentiation）记录运算及其局部导数，在数值计算中执行这个过程，不是用符号软件给整个网络生成一个巨大解析式，也不是用有限差分逐个扰动全部参数。

冻结 $W_0$ 表示不对它进行参数更新，通常也不保存其参数梯度，但计算 $\nabla_x\mathcal L=W_0^{\mathsf T}g$ 仍可能必不可少。如果 $x$ 依赖更早的可训练适配器，就必须让梯度穿过这个冻结层。只有当一个子图的所有输入与参数都不需要梯度，且不存在更早可训练对象依赖这条路径时，才可安全地把它完全排除在反向图之外。参数的可训练标记和整个运算的梯度记录开关属于不同层次，框架行为应结合自动微分文档检查\cite{pytorch-autograd}。

\subsection{随机梯度下降与小批量}

设 $\phi$ 汇集所有可训练参数。随机梯度下降（stochastic gradient descent, SGD）在第 $t$ 步用小批量 $\mathcal B_t$ 估计梯度，进行
$\phi_{t+1}=\phi_t-\eta_tg_t$，其中 $\eta_t>0$ 是学习率。若对有限训练集均匀采样，并正确平均单样本梯度，则在给定当前参数时 $g_t$ 是经验风险梯度的无偏估计；若使用不均匀采样、不同长度归一化或有状态的数据处理，就需要重新检查这个性质。

梯度累积（gradient accumulation）把多个微批次的梯度加起来，再执行一次优化器更新。若希望等价于一个按有效词元平均的大批次，应累积各微批次损失的分子，并按总有效词元数缩放。简单地将每个微批次的平均损失再除以微批次数，只在各微批次有效词元数相同时实现同样的加权。随机 dropout、有限精度和不同计算顺序还可能造成额外数值差异。

\subsection{Adam、权重衰减与状态开销}

自适应矩估计（adaptive moment estimation, Adam）为每个可训练标量维护梯度的一阶、二阶矩估计。设 $g_t$ 是第 $t$ 步梯度，
\[
m_t=\beta_1m_{t-1}+(1-\beta_1)g_t,\qquad
v_t=\beta_2v_{t-1}+(1-\beta_2)g_t^{\odot2},
\]
其中 $0\leq\beta_1,\beta_2<1$，$\odot2$ 表示逐元素平方，初始状态取零。令
$\widehat m_t=m_t/(1-\beta_1^t)$、
$\widehat v_t=v_t/(1-\beta_2^t)$，则更新包含
$-\eta_t\widehat m_t/(\sqrt{\widehat v_t}+\varepsilon)$，所有除法与平方根逐元素进行，$\varepsilon>0$ 用于数值稳定。AdamW 另外施加解耦权重衰减，例如先乘 $1-\eta_t\lambda$ 再加入优化步，其中 $\lambda\geq0$。

在普通 SGD 下，对目标加入 $\lambda\|\phi\|_2^2/2$ 与相应权重衰减可得到相同的一步形式；在自适应预条件的优化器中，把正则梯度加入 $g_t$ 与解耦衰减通常不同。对于 LoRA，衰减两个因子也不等于直接衰减乘积增量，后文将进一步分析。因此描述实验不能只写“使用正则化”，还应说明正则对象、系数和优化器实现。

梯度裁剪（gradient clipping）则限制一次更新所依据的梯度尺度，例如将总梯度乘上
$\min\{1,c/\|g_t\|_2\}$，其中 $c>0$ 是阈值，零梯度时保持不变。它有助于处理异常大梯度，但会改变更新方向的长度，不能替代检查标签错位、学习率过大和非有限损失。学习率预热、衰减和提前停止同样属于算法配置，需要在对照实验中保持清楚。

\section{理解插入位置：Transformer 的必要结构}
\label{sec:transformer}

\subsection{从词元嵌入到注意力}

Transformer 的基本结构见文献\cite{vaswani2017}。为统一本文的列向量约定，将长度为 $T$ 的序列表示为
$H=[h_1,\ldots,h_T]\in\mathbb R^{d\times T}$，每列对应一个位置的隐藏向量。词元嵌入（token embedding）把词表索引映射为向量，位置相关机制提供顺序信息。词表大小、隐藏维度和序列长度是不同的量：增加序列长度不改变某个线性投影的参数形状，却会增加激活与注意力计算。

以单个注意力头为例，令
$W_Q,W_K\in\mathbb R^{d_k\times d}$、
$W_V\in\mathbb R^{d_v\times d}$，得到
$Q=W_QH$、$K=W_KH$、$V=W_VH$。查询（query）与键（key）决定位置之间的匹配程度，值（value）提供被汇聚的内容。定义行归一化的注意力矩阵
\[
P=\operatorname{softmax}_{\mathrm{row}}
\left(\frac{Q^{\mathsf T}K}{\sqrt{d_k}}+M_{\mathrm{attn}}\right)
\in\mathbb R^{T\times T},
\qquad O=VP^{\mathsf T}.
\]
其中 $P_{ij}$ 表示查询位置 $i$ 对键位置 $j$ 的权重，$M_{\mathrm{attn}}$ 是注意力掩码，禁止访问的位置在理想表达中置为负无穷。实际数值实现需避免某一整行全部被屏蔽，否则 softmax 可能没有有效定义。

分母 $\sqrt{d_k}$ 的一个动机是控制点积尺度：若查询与键各分量相互独立、均值为零、方差为一，则点积的方差为 $d_k$，缩放后为一。真实隐藏特征一般不严格满足独立假设，所以这是初始化尺度的解释，而非对任意训练状态的精确分布结论。多头注意力将多个头的输出拼接，再通过输出投影 $W_O$ 变换回模型维度。

\subsection{注意力投影与前馈网络各自控制什么}

若对 $W_Q$ 或 $W_K$ 加入 LoRA，首先改变的是查询与键，从而影响注意力分数及其归一化结果；若对 $W_V$ 加入 LoRA，则直接改变被汇聚的内容；若对 $W_O$ 加入 LoRA，则改变各头信息混合后的输出表示。这样的功能区分有助于理解插入位置，但网络后续层会相互耦合，不能据此断言某一个矩阵只负责某一种语义能力。

前馈网络（feed-forward network, FFN）对每个位置独立施加非线性变换。例如
$\operatorname{FFN}(h)=W_2\sigma(W_1h+b_1)+b_2$，其中中间宽度为 $d_{\mathrm{ff}}$，
$W_1\in\mathbb R^{d_{\mathrm{ff}}\times d}$、
$W_2\in\mathbb R^{d\times d_{\mathrm{ff}}}$。门控结构还可能具有第三个投影矩阵。对这些矩阵适配可以改变特征组合与非线性处理方式，因此只适配注意力与同时适配前馈层是在选择不同的可训练函数族。

残差连接（residual connection）把一个分支的输出与输入相加，归一化层则调整特征尺度。LoRA 自身的加法形式与残差思想相似，但它通常发生在某个线性权重的作用内部。新增分支只要保持为线性形式，就能合并进该权重；整个 Transformer 仍有 softmax、激活函数和归一化等非线性，不能因为某一增量可合并就把网络看成一个低秩线性模型。

\subsection{不同模型实现中的形状差异}

一些实现把查询、键和值的权重拼成一个大矩阵，一些实现使用不同数量的查询头与键值头。此时矩阵可能不是 $d\times d$，参数量必须按实际形状计算。卷积、嵌入层和共享权重也可以设计低秩适配，但矩阵化方式、权重绑定和合并逻辑需要另外核查。模块名字相似不保证存储方向相同，本文数学约定是 $W\in\mathbb R^{m\times n}$、$h=Wx$，程序常用行批次 $X_{\mathrm{row}}W^{\mathsf T}$。

因此，开始训练前应检查目标模块的类型、输入输出维度和共享关系，并打印实际可训练参数数量。只看到配置中写了一个模块名称，并不能证明所有预期层都已注入适配器。若模型使用融合投影、特殊张量并行布局或自定义算子，现成注入规则可能只覆盖部分权重，甚至覆盖不到任何权重。

\section{LoRA 参数化、参数预算与前向计算}
\label{sec:lora-definition}

\subsection{完整的基本定义}

对选中的冻结权重 $W_0\in\mathbb R^{m\times n}$，给定整数
$1\leq r\leq\min(m,n)$、固定系数 $\alpha>0$，定义
\begin{equation}
W_{\mathrm{eff}}=W_0+\Delta W,\qquad
\Delta W=sBA,\qquad
s=\frac{\alpha}{r},
\quad A\in\mathbb R^{r\times n},\ B\in\mathbb R^{m\times r}.
\label{eq:lora}
\end{equation}
这里 $W_{\mathrm{eff}}$ 是实际生效权重，$A,B$ 是训练变量，$s$ 是给定缩放。该参数化对应标准 LoRA 的基本形式\cite{hu2021}。具体软件还可能支持其他缩放、初始化、偏置训练或变体；这些配置一旦改变，就应重新说明实际训练的参数集合。

对列批次 $X\in\mathbb R^{n\times b}$，其中 $b$ 是展平后的样本或词元数，前向计算为
\begin{equation}
H=W_0X+sB(AX).
\label{eq:lora-forward}
\end{equation}
训练时不需要先形成完整的 $BA$ 再与 $X$ 相乘。先算 $AX$，再算 $B(AX)$，既保持同一数学结果，又让新增分支的计算依赖于较小的中间维度。偏置若存在，则向每列输出添加同一个偏置向量。

\subsection{多层模型中的约束问题}

设被适配层的索引集合为 $\mathcal S$，第 $\ell$ 层尺寸为
$m_\ell\times n_\ell$，秩上限为 $r_\ell$，所有可训练因子汇集为
$\phi=\{A_\ell,B_\ell:\ell\in\mathcal S\}$。训练问题是
\[
\min_{\phi}\widehat R\left(
\{W_{\ell,0}+s_\ell B_\ell A_\ell\}_{\ell\in\mathcal S},
\{\text{其余冻结参数}\}\right).
\]
从表达集合上看，它等价于只允许选中层的增量满足
$\operatorname{rank}(\Delta W_\ell)\leq r_\ell$，因为非零缩放可吸收进因子。这里的等价仅指可表示的权重集合：在因子空间中执行梯度下降，不等于每步先做无约束梯度更新，再投影到低秩集合。

总适配参数量为
\begin{equation}
P_{\mathrm{LoRA}}=\sum_{\ell\in\mathcal S}
r_\ell(m_\ell+n_\ell).
\label{eq:parameter-count}
\end{equation}
若额外训练偏置、分类头、嵌入或幅度向量，还应把它们计入。比例的分母也须说明：相对于全部基座参数的比例，和相对于被适配矩阵原始参数的比例可能差很多。

\begin{example}[同样秩下的单层与整网预算]
假设有 $32$ 个结构相同的块，每块只适配两个 $4096\times4096$ 投影，均取 $r=8$。不计偏置时，总适配参数为
$32\cdot2\cdot8\cdot8192=4194304$。被适配的原矩阵共有
$32\cdot2\cdot4096^2=1073741824$ 个元素，两者之比为 $1/256$。这不是相对某个完整大模型的比例，因为这里没有指定其嵌入、前馈网络及其他参数的数量。
\end{example}

\subsection{低秩分支中的随机失活}

随机失活（dropout）可以施加在低秩分支输入上。设失活概率 $0\leq p<1$，掩码各分量独立服从成功概率 $1-p$ 的伯努利分布（Bernoulli distribution），记
$D_p(x)=m\odot x/(1-p)$，则训练时采用
$h=W_0x+sBA D_p(x)$。在给定 $x$ 时，
$\mathbb E[D_p(x)]=x$，所以该线性层的分支输出期望等于不失活的输出。

然而，后续网络存在非线性，一般有
$\mathbb E[f(D_p(x))]\neq f(\mathbb E[D_p(x)])$。不能把测试时关闭失活解释为对整个随机训练网络输出的精确求平均。失活是一种正则化机制，既可能降低过拟合，也可能在数据不足以支持高噪声训练时妨碍优化，其作用需要通过验证集检验。推理合并讨论默认失活已经关闭。

\section{双因子的梯度、初始化与手算训练}
\label{sec:lora-gradient}

\subsection{从有效权重梯度推导因子梯度}

先假设没有失活，损失可微，且对有效权重的梯度为
$G=\nabla_{W_{\mathrm{eff}}}\mathcal L\in\mathbb R^{m\times n}$。由于 $W_0$ 与 $s$ 固定，
$\mathrm dW_{\mathrm{eff}}=s(\mathrm dB)A+sB(\mathrm dA)$。使用迹的循环性质，
\[
\mathrm d\mathcal L
=s\operatorname{tr}(G^{\mathsf T}(\mathrm dB)A)
+s\operatorname{tr}(G^{\mathsf T}B(\mathrm dA)).
\]
把它分别整理为因子与其微分的 Frobenius 内积，得到
\begin{equation}
\nabla_B\mathcal L=sGA^{\mathsf T},\qquad
\nabla_A\mathcal L=sB^{\mathsf T}G.
\label{eq:factor-gradients}
\end{equation}
形状检查是推导的一部分：第一式为 $(m\times n)(n\times r)=m\times r$，第二式为 $(r\times m)(m\times n)=r\times n$。

若从批次计算图出发，记 $Z=AX$、$E=\nabla_H\mathcal L$，则
\[
\nabla_B\mathcal L=sEZ^{\mathsf T},\qquad
\nabla_A\mathcal L=sB^{\mathsf T}EX^{\mathsf T},\qquad
\nabla_X\mathcal L=W_0^{\mathsf T}E+sA^{\mathsf T}B^{\mathsf T}E.
\]
其中 $E$ 已包含损失归一化的系数，不能再无依据地除一次批量大小。训练时实际可沿小矩阵路径计算因子梯度，不必存储完整 $G$；数学上用 $G$ 只是为了简化表达。若启用输入失活，则因子梯度中的 $X$ 应改为失活后的输入，并在输入梯度上应用对应掩码。

\subsection{一个因子为零为什么仍能开始学习}

常用初始化令 $A_0$ 随机非零、$B_0=0$。于是 $\Delta W_0=0$，初始模型与预训练模型保持相同的前向函数，同时
\[
\left.\nabla_A\mathcal L\right|_0=0,\qquad
\left.\nabla_B\mathcal L\right|_0=sG_0A_0^{\mathsf T}.
\]
右侧通常非零，所以第一步先更新 $B$，后续 $B$ 非零后 $A$ 也能得到任务梯度。这里“通常”不是“必然”：若当前任务梯度为零，或者恰好与初始化子空间正交，第一步也可能没有变化。官方实现可用于核对具体初始化与存储方向\cite{lora-code}，不同代码版本采用的随机分布未必完全相同。

如果把 $A_0$、$B_0$ 同时置零，则式\eqref{eq:factor-gradients}中的两个任务梯度都为零。对没有额外扰动、没有非零优化器初始状态的标准一阶方法，两个因子会一直停在零点；常见的二次正则与权重衰减也不会把零变成非零。反过来令 $A_0=0$、$B_0$ 随机也能保持初始增量为零，但最先学习的因子与初期梯度子空间会交换，训练轨迹一般并不相同。

\subsection{一个完整的一步更新}

令 $W_0=I_2$，$r=1$，$s=1$，
$A_0=(1,0)$、$B_0=(0,0)^{\mathsf T}$，给定
$x=(2,1)^{\mathsf T}$、目标 $y=(0,1)^{\mathsf T}$，采用平方损失
$\mathcal L=\frac12\|h-y\|_2^2$。初始输出为 $h_0=(2,1)^{\mathsf T}$，损失为 $2$，上游梯度为 $g_0=(2,0)^{\mathsf T}$，中间特征为 $A_0x=2$。因此
\[
\nabla_B\mathcal L=(4,0)^{\mathsf T},\qquad
\nabla_A\mathcal L=(0,0).
\]
用学习率 $\eta=0.1$ 同时更新两个因子，得到
$B_1=(-0.4,0)^{\mathsf T}$、$A_1=(1,0)$。新输出为
$h_1=(1.2,1)^{\mathsf T}$，新损失为 $0.72$。

在这个新状态，$g_1=(1.2,0)^{\mathsf T}$，于是
$\nabla_A\mathcal L=B_1^{\mathsf T}g_1x^{\mathsf T}
=(-0.96,-0.48)$，第二个因子也开始通过任务梯度调整。该例既展示了零初始化的作用，也说明梯度必须由同一个前向状态计算。若先更新 $B$ 再用更新后的 $B$ 计算同一步的 $A$ 梯度，就已经换成另一种交替算法。

\subsection{随机初始化决定早期的可见方向}

在 $B_0=0$ 的第一步 SGD 中，$B_1=-\eta sG_0A_0^{\mathsf T}$，因而
\begin{equation}
\Delta W_1=-\eta s^2G_0A_0^{\mathsf T}A_0.
\label{eq:first-step}
\end{equation}
矩阵 $A_0^{\mathsf T}A_0$ 的秩至多为 $r$，所以初期更新会经过初始化给定的输入方向。若 $A_0$ 的行恰好正交且单位化，它才是正交投影矩阵；一般随机初始化并不满足这一条件，因此不能无条件称它为正交投影。

假设 $A_0$ 的各元素独立、均值为零、方差为 $\tau^2$，且 $G_0$ 在零增量处不依赖 $A_0$，则
$\mathbb E[A_0^{\mathsf T}A_0]=r\tau^2I_n$，从而初始化平均意义下
$\mathbb E[\Delta W_1]=-\eta s^2r\tau^2G_0$。单次运行仍然是低秩更新，而多个随机初始化的平均可以是满秩矩阵，这没有矛盾。该式还说明秩、随机尺度和缩放共同影响有效初始步长。

\section{优化几何：为什么训练因子不同于直接训练权重}
\label{sec:geometry}

\subsection{乘积参数化改变了更新方向}

假设在某一步使用同一学习率 $\eta$ 的普通 SGD，不含失活、正则项和动量，令
$A^+=A-\eta sB^{\mathsf T}G$、
$B^+=B-\eta sGA^{\mathsf T}$。直接展开新乘积可得精确恒等式
\begin{equation}
\Delta W^+-\Delta W
=-\eta s^2\left(GA^{\mathsf T}A+BB^{\mathsf T}G\right)
+\eta^2s^3GA^{\mathsf T}B^{\mathsf T}G.
\label{eq:effective-step}
\end{equation}
第一阶项既对输入方向进行重加权，也对输出方向进行重加权；第二阶项来自两个因子同时变化。因此因子训练相当于引入了随状态变化的优化几何，而不是简单地从全参数梯度中抽出几个元素。

这也解释了为什么不能照搬全参数微调的学习率。即使两个方法的当前有效权重相同，同样的标量学习率也可能产生不同大小、不同方向的有效更新。使用 Adam 后，逐元素的自适应缩放又进一步依赖因子坐标，式\eqref{eq:effective-step}不再直接描述实际更新；该式的用途是揭示机制，而不是给所有优化器提供统一动力学公式。

\subsection{非唯一分解与尺度不平衡}

对非零标量 $c$，$B'=cB$、$A'=A/c$ 表示相同增量，但
$\nabla_{B'}\mathcal L=(1/c)\nabla_B\mathcal L$、
$\nabla_{A'}\mathcal L=c\nabla_A\mathcal L$。因此同一个有效模型可能因为因子尺度不同而呈现不同梯度尺度。过大的尺度失衡会使某个因子更新很快、另一个更新很慢，并可能增加有限精度下的溢出或舍入问题。

这不是说两个因子的范数必须完全相等，而是说明监控 $\|A\|_F$、$\|B\|_F$ 和 $\|\Delta W\|_F$ 提供不同信息。某个因子范数增加而有效更新几乎不变，可能只是分解坐标发生变化。比较两次训练时，直接比较因子元素或范数未必公平；比较有效增量、输出变化和任务指标通常更有意义。

\subsection{因子正则化与核范数的联系}

核范数（nuclear norm）定义为矩阵奇异值之和，
$\|D\|_*=\sum_j\sigma_j(D)$。当 $D$ 的秩不超过 $r$ 时，有
\begin{equation}
\min_{BA=D}\frac12\left(\|A\|_F^2+\|B\|_F^2\right)=\|D\|_*,
\label{eq:nuclear}
\end{equation}
其中最小化在尺寸为 $r\times n$ 与 $m\times r$ 的因子上进行。该关系说明，若对所有表示同一乘积的分解充分优化，对因子的二次惩罚可以与鼓励谱集中联系起来。

\begin{proof}
写 $BA=\sum_kb_ka_k^{\mathsf T}$。核范数的三角不等式及秩一矩阵的奇异值公式给出
$\|BA\|_*\leq\sum_k\|b_k\|_2\|a_k\|_2
\leq\frac12(\|B\|_F^2+\|A\|_F^2)$。
若 $D=U_D\Sigma_DV_D^{\mathsf T}$ 为仅保留非零奇异值的分解，取
$B=U_D\Sigma_D^{1/2}$、
$A=\Sigma_D^{1/2}V_D^{\mathsf T}$，再按需补零，则达到等号。这里使用了核范数作为范数的三角不等式，不重证该范数性质。
\end{proof}

对于实际增量 $\Delta W=sBA$ 且 $s>0$，右侧对应
$\|\Delta W\|_*/s$。不过，有限步的 AdamW 因子训练并不自动等价于求解一个核范数正则化的权重问题：训练没有每步在等价分解中求最小值，优化器也未必对应把二次惩罚加入目标。应把式\eqref{eq:nuclear}视为明确优化条件下的数学联系，而非对所有实现的解释定理。

\subsection{低秩限制如何穿过多层网络}

在仅有两个线性层的例子中，总映射为
$(W_{2,0}+D_2)(W_{1,0}+D_1)$，相对基座的变化是
$W_{2,0}D_1+D_2W_{1,0}+D_2D_1$。即使各 $D_\ell$ 都是秩一，不同项相加也可能产生多个独立方向。深层网络还有输入相关的非线性门控，因此“每层小秩”不意味着整个模型在所有输入上的函数变化具有同一个小秩。

在固定输入附近，若模型关于参数可微，可线性化为
$f_{\theta_0+\delta}(x)\approx f_{\theta_0}(x)+J_{\theta_0}(x)\delta$，其中 $J_{\theta_0}(x)$ 是输出对参数的雅可比矩阵（Jacobian matrix）。LoRA 限制可实现的 $\delta$，而真正影响输出的是 Jacobian 映射后的方向。不同层相同的秩预算可能产生很不相同的函数变化，这为按层选择模块与分配预算提供了理由，但不直接决定最优方案。

\section{秩、缩放、初始化与表达能力}
\label{sec:rank-scaling}

\subsection{增大秩能保证什么}

若目标模块相同、缩放非零且不额外限制因子范数，把秩从 $r_1$ 增加到 $r_2\geq r_1$ 后，可通过补零表示原来的任何增量。因此较大秩的可表示集合包含较小秩的集合，理想全局最小训练损失不会因秩增加而变差。这个结论讨论的是不加不同正则项的全局最优值，不是某次有限时间训练实际达到的损失。

验证误差也没有相同的单调性。更大的秩增加拟合自由度，可能改善欠拟合，也可能拟合标签噪声；它还会改变初始化统计、缩放和优化难度。对某些任务，增加目标模块覆盖率比提高少数模块的秩更有效；对另一些任务，关键层需要更充分的方向。只有比较受控实验才能分辨这些情形。

\subsection{为什么缩放与秩不能分开看}

将式\eqref{eq:first-step}与随机初始化分析结合，可见在固定 $\alpha$、$\tau$ 和 SGD 学习率时，标准缩放 $s=\alpha/r$ 使初始化平均步长系数为
$\eta\alpha^2\tau^2/r$。若采用 $s=\alpha/\sqrt r$，该系数变为
$\eta\alpha^2\tau^2$。秩稳定 LoRA（rank-stabilized LoRA, rsLoRA）研究了后一种缩放\cite{kalajdzievski2023}。这里的一步期望推导依赖前文的独立性与初始化条件，不能替代该方法对训练过程的完整分析。

另一种直觉来自外积和：若每个方向对某个输出分量提供近似独立、零均值、相近方差的随机贡献，则 $r$ 项相加的方差与 $r$ 成比例，乘 $1/\sqrt r$ 后方差保持同一量级。训练后的方向通常会相关，且零初始化时分支输出本来就是零，所以这只是解释缩放选择的辅助模型，不能声称它在所有训练阶段精确成立。

若同时改变 $\alpha$、学习率、秩和初始化方差，就很难判断性能变化来自哪一个因素。即使保持 $s$ 不变，参数个数与优化轨迹仍会改变；即使保持某种初始有效步长不变，后续自适应优化也可能不同。严谨的秩消融至少应记录全部这些配置，而不只在表格中列出 $r$。

\subsection{低秩适配为什么可能有效，又会在哪里失败}

若任务只要求在已有表示上重新组合少量特征，低秩增量可能足够表达有用变化；冻结其余权重还能约束训练自由度，使小数据下的过拟合更容易控制。这是可检验的机制假设，并不意味着模型把一切知识都预先存储完毕，也不意味着 LoRA 只能改变风格而不能学习新关联。任务所需变化、数据覆盖与优化能力共同决定结果。

可以构造一个无法绕过的表达限制：在单个线性层中，若目标增量恰为 $I_n$，输入二阶矩为 $I_n$，而 $r<n$，则由定理\ref{thm:truncated-svd}，任意秩不超过 $r$ 的增量都满足输出平方误差至少为 $n-r$。如果模型只有这一层，增加训练步数不能消除该下界。深层网络可以通过多个位置分担变化，因此不能把这个单层下界直接当作真实语言模型的误差下界，但它足以反驳“低秩没有表达损失”的绝对说法。

还应区分矩阵秩与任务的内在维数（intrinsic dimension）。前者是某个线性映射的代数性质，后者可能指一类解、参数变化或数据结构所需的有效坐标数量。二者之间可能有关联，却没有无条件的一一对应。不能从低维数据流形直接推出每一层的最优权重增量都具有相同的小秩。

\section{参数、显存与计算成本的逐项分析}
\label{sec:resources}

\subsection{训练显存由哪些部分组成}

设基座共有 $P_0$ 个参数，可训练新增参数共有 $P_a$ 个。训练峰值显存可以概念性地拆成权重存储、梯度存储、优化器状态、保留激活和临时工作区。前几项较容易按参数数量估算，后两项依赖序列长度、微批次、层数、算子实现与重计算策略。设备显存分配器还可能保留缓存，因此软件报告的已分配与已保留显存也不完全相同。

作为明确假设下的估算，设全参数训练为每个参数存储两字节计算权重、两字节梯度、四字节主权重和两个四字节 Adam 状态，共十六字节，则参数相关显存约为 $16P_0$ 字节。若基座冻结，以两字节存储，每个适配参数仍使用上述十六字节，则参数相关显存约为
$2P_0+16P_a$ 字节。不同混合精度方案不一定存在独立主权重，梯度和状态也可能采用其他精度，所以这些系数不是通用常数。

例如在这一假设下，$P_0=7\times10^9$、$P_a=10^7$ 时，全参数训练的参数相关存储约为 $112$ GB，而冻结基座加适配参数约为 $14.16$ GB。这里 GB 按 $10^9$ 字节计，尚未包括激活和工作区，也未加入分布式分片；因此不能据此宣布某一显存容量必然可以完成训练。特别是长序列任务，激活可能成为主要瓶颈。

\subsection{为什么冻结权重仍然需要大量计算}

对 $n\times b$ 的批次输入，基座前向乘法的乘加次数量级为 $mnb$，低秩分支增加约 $rb(n+m)$ 次乘加。这里把一次乘加作为一个计数单位；若按浮点运算次数统计，一次乘法和一次加法常计为两个操作，报告时要统一口径。低秩分支的参数少，不代表基座前向计算会消失。

全参数线性层的反向计算通常包含权重梯度与输入梯度两项，两者规模都与前向乘法相近。冻结基座可省去它的权重梯度，但为了训练更早的适配器，往往仍需计算输入梯度。只看一个理想稠密线性层，可以把三项主要矩阵乘法减少到两项，再加上适配分支成本；这只是算术模型，不是整体训练吞吐的固定比例。

真实速度还取决于矩阵规模是否适合硬件、内存带宽、内核启动、通信、数据加载和量化解码。很小的低秩矩阵乘法可能难以充分利用设备；省下显存则可能允许更大微批次，间接提升吞吐。应分别记录每秒有效词元数、总训练时间和峰值显存，不能只看适配器文件大小来判断训练效率。

\subsection{激活重计算、混合精度与分布式训练}

激活检查点或激活重计算（activation checkpointing）只保留部分中间结果，反向时重新执行一些前向运算，以计算换显存。它与 LoRA 解决的是不同存储部分，可以组合使用。重计算需要保持与原前向一致的随机行为；若失活掩码被错误重采样，得到的梯度可能不对应原先的损失计算。

混合精度（mixed precision）让不同操作使用不同浮点格式。较低精度节省带宽与存储，较高精度适合累积、归一化或敏感的统计量。即使基座冻结，低秩分支的梯度仍可能因范围过小而下溢，或因尺度过大而溢出。发现非有限梯度时，应先定位哪个张量开始异常，再检查精度、损失归一化、学习率和初始化，而不是一律增大秩。

在数据并行中，只同步可训练适配参数的梯度可减少这部分通信，但每个设备仍可能保存完整基座。参数分片、优化器状态分片与流水并行具有不同的通信和存储取舍。报告多设备结果时，应说明显存是单设备峰值还是总和，吞吐是单设备还是全局，适配器保存是否已经聚合。参数高效方法可以与分布式系统组合，但不能替代明确的系统配置。

\section{量化前置知识与 QLoRA 的工作机制}
\label{sec:quantization}

\subsection{量化与低秩分解解决不同问题}

量化（quantization）用有限集合中的值或编码近似原数值，主要减少每个权重的存储位数；低秩分解减少表示某个矩阵所需的独立参数数量。一个矩阵可以低位但满秩，也可以高精度但低秩。把基座量化与训练低秩增量结合，是同时减少基座存储与可训练状态的一种方式，不应把二者视为同一个压缩步骤。

最简单的对称均匀量化可写为
$q=\operatorname{clip}(\operatorname{round}(w/c),q_{\min},q_{\max})$、
$\widehat w=cq$，其中 $c>0$ 是尺度，$q$ 是整数编码。对于没有截断的数值，最近舍入的绝对误差至多为 $c/2$；超出编码范围的数值还会产生截断误差。较大的尺度覆盖范围更广但分辨率更粗，较小的尺度分辨率更细但更容易截断异常值。

分块量化（blockwise quantization）为每个小块单独保存尺度，避免少数大值使整个张量的量化过粗。然而尺度本身也占存储，所以“每权重四位”通常仅描述主体编码，不代表整个文件精确地只占参数数目的一半字节。例如每 $64$ 个四位编码另外保存一个三十二位尺度，忽略其他信息后，每个参数平均占 $4+32/64=4.5$ 位。

\subsection{量化基座上的可微计算}

记冻结的量化编码为 $q_0$，尺度等元数据为 $c_0$，反量化算子为
$\mathcal D$，则实际参与运算的近似基座为
$\widehat W_0=\mathcal D(q_0,c_0)$。适配层执行
\begin{equation}
h=\widehat W_0x+sBAx.
\label{eq:qlora-forward}
\end{equation}
编码和尺度在这里不训练，因子梯度仍由普通链式法则计算；输入梯度包含
$\widehat W_0^{\mathsf T}g$。因此无需对离散舍入函数求导，就能把梯度传回适配器。这与需要学习量化前权重、因而可能使用梯度近似的量化感知训练（quantization-aware training, QAT），是不同的训练安排。

储存时使用低位编码，并不意味着矩阵乘法、激活和梯度都在同样位宽下进行。实现可以按块反量化并与计算融合，而不必在显存中永久保存完整高精度副本。计算精度、累积精度和存储精度应分别说明，不能把“四位训练”理解为所有中间量只含十六个可取值。

\subsection{QLoRA 的三个关键部件}

QLoRA 在冻结量化基座上训练低秩适配器，并讨论了 NF4、双重量化和分页优化器三项机制\cite{dettmers2023}。NF4 即四位 NormalFloat，使用适合近似正态权重分布的非均匀码本；双重量化进一步量化分块尺度等量化常数，降低元数据开销；分页优化器借助内存分页处理状态存储的峰值需求。这些机制分别涉及编码、元数据和内存管理，不是三个不同的低秩分解公式。

非均匀码本可以在数值更密集的区域放置更多代表点。它是否适合某一权重分布，与异常值、分块归一化和误差指标有关，不能宣称对任意分布都最优。分页可以缓解显存峰值，但数据在不同存储层之间迁移会带来延迟，其收益依赖实际访问与硬件环境。本文不把论文中特定配置的性能或容量结论推广为所有模型的保证。

\subsection{量化误差会怎样影响任务适配}

设 $\widehat W_0=W_0+E_q$，则 QLoRA 的实际有效权重为
$W_0+E_q+\Delta W$。训练因子既要完成目标任务所需变化，也可能部分补偿量化误差，但秩限制意味着不能一般性地精确抵消任意 $E_q$。如果量化误差主要落在任务不敏感的方向，影响可能较小；如果落在敏感方向，则可能需要更高精度、调整分块方式或改变目标模块。

从线性输出看，量化误差造成的均方影响为
$\operatorname{tr}(E_qC_XE_q^{\mathsf T})$，与式\eqref{eq:data-weighted-error}相同。由此可见，单纯报告权重均方量化误差不足以描述任务退化；数据分布和下游敏感性也重要。实际比较应包括同一基座的高精度 LoRA 与量化 LoRA，并保持训练数据、词元预算和验证规则一致，才能分辨量化与其他配置的影响。

\section{从公式到实现：一个可核查的训练流程}
\label{sec:implementation}

\subsection{模块实现的最小结构}

下面代码展示 PyTorch 中单个线性层的教学实现，使用本文的 $W_0\in\mathbb R^{m\times n}$ 约定。框架的线性算子接受末维为输入特征的张量，权重形状及运算定义见文档\cite{pytorch-linear}，所以实际批次运算与式\eqref{eq:lora-forward}互为转置表达。代码刻意只包含冻结基座、两个可训练因子和缩放，未实现量化、失活、分布式切分与自动合并；这些功能会引入额外状态，适合在验证基础关系后逐项加入。

\begin{verbatim}
import math
import torch
from torch import nn
from torch.nn import functional as F

class LoRALinear(nn.Module):
    def __init__(self, base, rank, alpha):
        super().__init__()
        if not isinstance(base, nn.Linear):
            raise TypeError("base must be nn.Linear")
        if not 1 <= rank <= min(base.in_features,
                                base.out_features):
            raise ValueError("invalid rank")
        if not math.isfinite(alpha) or alpha <= 0:
            raise ValueError("alpha must be positive")
        self.base = base
        self.base.requires_grad_(False)
        self.scale = alpha / rank
        opts = dict(device=base.weight.device,
                    dtype=base.weight.dtype)
        self.A = nn.Parameter(torch.empty(
            rank, base.in_features, **opts))
        self.B = nn.Parameter(torch.zeros(
            base.out_features, rank, **opts))
        nn.init.normal_(self.A, std=1 / math.sqrt(
            base.in_features))

    def forward(self, x):
        low_rank = F.linear(F.linear(x, self.A), self.B)
        return self.base(x) + self.scale * low_rank

    @torch.no_grad()
    def merged_weight(self):
        return (self.base.weight
                + self.scale * (self.B @ self.A))
\end{verbatim}

这个类直接持有传入的基座层，并将其冻结；若其他模块也共享该对象，其参数的冻结状态也会改变。构造时要求传入普通浮点线性层，不能把它直接当作任意量化层的包装器。返回合并权重的方法只生成一个张量，不修改原层，因此不会因重复调用而把增量重复加到基座上。若基座具有偏置，部署为普通线性层时还需复制原偏置。

\subsection{应当验证哪些不变量}

初始化后，在相同输入和计算精度下，适配层输出应与基座输出一致，因为 $B=0$。执行反向传播后，基座参数不应得到新的梯度，可训练因子的梯度形状应与参数相同；在普通任务损失下，第一步 $A$ 的梯度应为零，$B$ 的梯度则应与手算结果一致。这里不能要求每个随机样本都使 $B$ 梯度严格非零，因为输入或上游梯度可能本来就是零。

训练一步后，应比较参数的实际变化，确认优化器参数组确实包含 $A,B$。仅看到损失具有梯度函数，不能证明目标因子被加入优化器；仅看到参数被标记为可训练，也不能证明损失依赖它。关闭随机失活后，再比较未合并输出与使用
$W_0+sBA$ 的普通线性输出，并根据计算精度设置误差容忍度。

有限差分可以在小型双精度例子中检查解析梯度。对某个因子元素施加正负扰动 $\varepsilon$，计算
$[\mathcal L(\phi+\varepsilon e_j)-\mathcal L(\phi-\varepsilon e_j)]/(2\varepsilon)$，并与自动微分结果比较。$\varepsilon$ 太大产生截断误差，太小产生舍入误差，应选中间尺度并检查若干不同扰动。大模型不适合逐元素做这种昂贵检查，小矩阵已经足以发现转置、缩放和归一化错误。

\subsection{训练循环的完整逻辑}

一个可复现流程需要先固定基座与分词器版本，再构造训练、验证和测试划分，明确对话模板、最大长度、特殊词元和监督掩码。随后选择目标模块、秩、缩放、初始化及额外训练参数，注入适配器，并统计实际参数量。优化器应在注入完成后构建，否则新参数可能遗漏，或者参数对象被替换后优化器仍持有旧对象。

每个优化步骤依次执行前向计算、按有效监督位置归一化损失、反向传播、必要的梯度累积与裁剪、优化器更新、学习率更新以及梯度清零。若采用混合精度缩放，梯度裁剪应作用于按所用框架要求还原尺度后的梯度。验证时切换到评估行为，关闭训练用失活，以独立验证指标选择检查点，并在最终确定配置后使用测试集。

终止条件可以是预定有效词元预算、最大更新步数，或验证指标在预定耐心范围内没有改善。提前停止规则应在实验前明确，避免看着测试集反复挑选停止点。最后保存适配器权重、完整配置、基座标识、分词器与模板信息；若希望中断后精确续训，还需要优化器、调度器、随机数与数据读取状态。只保存推理所需的适配器文件，并不等于保存完整训练检查点。

\section{权重合并、多适配器组合与部署}
\label{sec:merging}

\subsection{什么时候可以精确合并}

在推理时，若 $A,B,s$ 固定且分支没有输入相关的非线性或随机操作，则
$W_0x+sBAx=(W_0+sBA)x$，所以可以预先形成
$W_{\mathrm{merged}}=W_0+sBA$。这是矩阵分配律给出的代数恒等式。合并后移除适配分支，单个线性层可恢复为一次普通矩阵乘法，因而没有额外分支运算；实际延迟仍应由部署后端测量。

若分支改为 $B\sigma(Ax)$，一般不存在一个固定矩阵与它对所有输入等价。若推理时仍动态选择适配器权重或保留输入相关门控，也不能一次性合并为所有请求共享的同一矩阵。多个任务共享同一基座时，保留适配器便于切换，分别保存合并模型则减少单请求分支计算但增加存储，二者对应不同的服务需求。

浮点运算不严格满足实数运算的结合律与分配律。先算 $B(Ax)$ 与先算 $(BA)x$ 的舍入顺序不同，所以合并前后应检查数值误差，而不要求所有位完全一致。合并的工程陷阱还包括重复加增量、把转置方向用反、遗漏缩放、同时保留已合并分支，以及从错误的基座版本开始。保留原始基座与合并状态标记比反复做加减更稳妥。

\subsection{量化后的合并不再是无误差重写}

量化模型实际训练的是 $\widehat W_0+\Delta W$。若先把它合并为高精度矩阵，再量化为
$\mathcal D(\mathcal Q(\widehat W_0+\Delta W))$，通常会产生新的舍入误差。因此重新量化后的部署模型不必与训练时模型一致。若直接把适配增量加到原始高精度 $W_0$ 而不是训练时的 $\widehat W_0$，又会改变基座，差值至少在单层权重意义上包含此前的量化误差。

实际交付应明确选择哪一种路径：保留原量化基座与独立适配器、在训练时反量化基座上合并，或在原高精度基座上合并后重新评估。不同路径可能各有用途，但不能把它们默认为数学上同一个模型。模型加载成功只验证格式兼容，不验证函数与质量保持不变。

\subsection{多个低秩更新如何相加}

设两个任务的更新为 $D_1=s_1B_1A_1$、$D_2=s_2B_2A_2$，希望形成
$D=\lambda_1D_1+\lambda_2D_2$，其中 $\lambda_1,\lambda_2$ 是给定实系数。可以精确写成
\[
D=
\begin{bmatrix}\lambda_1s_1B_1&\lambda_2s_2B_2\end{bmatrix}
\begin{bmatrix}A_1\\A_2\end{bmatrix},
\qquad
\operatorname{rank}(D)\leq r_1+r_2.
\]
实际秩可能更小，因为两个更新可能共享方向或相互抵消。如果部署要求仍保持原来的小秩，可对合成更新做低秩近似，但由第\ref{sec:svd}节可知，权重误差小不保证任务指标保持。

直接平均 $A$ 再平均 $B$ 一般不等于平均更新，因为
$(B_1+B_2)(A_1+A_2)$ 多出 $B_1A_2+B_2A_1$ 两个交叉项；因子还有基变换冗余，使按元素平均更加缺少不变性。应先定义想要组合的实际增量或函数，再选择实现方式。即使参数组合数学上正确，不同任务的梯度与行为也可能冲突，组合后的能力不是各适配器能力的自动并集。

\subsection{连续训练与任务干扰}

在任务一训练得到适配器后继续用同一个适配器训练任务二，优化器只看到当前数据时可能损害任务一表现。冻结基座限制了权重变化的范围，却不保证旧任务输出保持，因为所有输入仍会经过新的适配分支。保留独立适配器可以恢复各自的任务模型，但如何路由请求、如何处理跨任务输入和如何控制多适配器存储，是另外的系统问题。

若把任务一增量合并进基座，再新增一个秩 $r$ 适配器训练任务二，累计增量的秩上界可随阶段增加；这与始终在同一个固定秩乘积中训练不同。比较连续学习方法时，应统计累计存储、原始基座是否保留，以及每个阶段是否改变了可表达集合，避免把更大的累计自由度误记为同等预算。

\section{代表性变体及其改变的数学对象}
\label{sec:variants}

\subsection{自适应预算与幅度方向分解}

自适应 LoRA（AdaLoRA）关注不同矩阵不应机械获得相同预算的问题，通过类似奇异值分解的参数化及重要性评分动态分配增量预算\cite{zhang2023}。从优化问题看，它同时决定“哪些方向怎样更新”和“各处保留多少方向”。重要性由训练中的代理量估计，不等于知道真实总体风险下每个方向的贡献；若任务分布变化，原有评分也可能不再适用。

权重分解低秩适配（weight-decomposed low-rank adaptation, DoRA）将权重的幅度与方向分开训练，并用低秩结构调整方向\cite{liu2024}。为说明其含义，沿用原论文按列归一化的约定，把第 $j$ 列写为
\[
w_j=\rho_j\frac{v_j}{\|v_j\|_2},
\qquad
v_j=w_{0,j}+s(BA)_{:,j},\qquad \|v_j\|_2>0.
\]
其中 $\rho_j$ 是单独训练的幅度变量，初始可取 $\|w_{0,j}\|_2$。具体实现若转置存储权重，归一化轴也须相应转换，零范数则需额外处理。

这类参数化使实际增量 $w_j-w_{0,j}$ 同时受幅度与归一化方向影响，整个实际增量不再简单等于一个秩不超过 $r$ 的 $BA$。例如仅对不同列作不同倍数缩放，差值就可能具有较高秩。因而“方向部分使用低秩”与“总增量严格低秩”必须区分；这种区分比只记住变体名称更能帮助理解其表达能力。

\subsection{缩放、初始化与学习率是不同改进轴}

第\ref{sec:rank-scaling}节的 rsLoRA 改变秩相关缩放；初始化方法可以改变起始子空间或起始分解；为两个因子设置不同学习率则改变乘积空间中的优化几何；QLoRA 主要改变基座存储与计算路径。这些变化通常彼此独立，因此一个实验可以同时采用其中多个部件，但必须明确组合后实际执行的公式。

若某个初始化希望让初始 $BA$ 非零，同时仍保持原模型函数不变，就需要在冻结基座中相应减去该初始增量，或者采用其他补偿结构。仅把两个随机因子都设为非零而不补偿，会从一开始扰动模型。这是恒等式层面的要求，不依赖具体初始化算法的名称，也不能仅靠“随机量均值为零”保证每次实际模型完全不变。

\subsection{与其他参数高效方法的关系}

只训练偏置限制的是参数坐标，只训练输出头限制的是层位置，软提示（soft prompt）学习输入侧的连续向量，瓶颈适配模块常在隐藏表示上添加小型非线性网络，而 LoRA 限制的是选定线性权重的增量结构。它们可以拥有相似参数量，却对应不同函数族和推理开销。参数相同不等于能力相同，也不等于计算成本相同。

比较方法时，应说明控制的是总训练参数、总训练时间、峰值显存还是最终存储。四种预算往往不能同时严格相等。一个方法在固定显存下更强，可能因为它允许使用更大基座；在固定基座下更强，则可能来自参数化、优化或数据差异。清楚界定比较条件，才能把经验结论放在正确的范围内。

\section{超参数选择、实验设计与故障诊断}
\label{sec:evaluation}

\subsection{先建立可解释的基线}

有意义的实验首先需要未微调基座的表现，确认任务到底缺少什么；然后可以加入只训练任务头或其他简单基线，判断复杂适配是否必要。在资源允许时，再加入经过合理调参的全参数微调。若只比较未经调参的全参数方案与精心搜索的 LoRA，就混入了搜索预算差异，不能据此断言某种参数化普遍更好。

数据划分需要防止内容泄漏。同一文档切出的片段、同一模板生成的近重复样本、同一用户的强相关记录，可能不适合随机分到训练与测试两边。应按任务泛化目标选择分组、时间或来源划分。数据越小，训练指标越容易受偶然样本影响，更需要报告多个随机种子的均值与波动，而不是只展示最佳种子。

\subsection{按问题而不是按流行配置选择超参数}

可先在有限搜索预算内比较几个适中秩和两类模块覆盖方案，并为各方案采用相近的调参机会。若训练与验证都差，先检查目标和实现，再考虑增加训练预算、模块覆盖或秩；若训练好而验证差，优先检查泄漏、噪声、分布差异与过拟合。提高秩不是通用修复，降低秩也不是自动正则化成功的证明。

学习率通常比精确挑选某个秩更直接影响训练是否稳定，但最佳值仍依赖基座、数据、缩放、有效批量和优化器。可以使用对数尺度的小范围候选搜索，结合梯度范数、有效更新范数和验证曲线判断。固定某个默认学习率然后比较很多秩，可能只是比较不同有效步长下谁更容易训练。

设第 $\ell$ 层基座与增量范数为
$\|W_{\ell,0}\|_F$、$\|\Delta W_\ell\|_F$，可监控相对更新尺度
$q_\ell=\|\Delta W_\ell\|_F/(\|W_{\ell,0}\|_F+\epsilon)$，其中 $\epsilon>0$ 防止分母为零。该比值能发现层间尺度异常，却不是通用质量指标：基座范数很大时，某个任务关键方向的变化仍可能被总量掩盖，因此还应查看输出与任务指标。

\subsection{损失不下降时的排查逻辑}

首先在少量样本上检查是否存在有效监督词元，标签是否与预测位置正确对齐，回答是否被最大长度截断，输入模板是否符合模型预训练或指令格式。然后检查实际可训练参数是否非空，梯度是否到达因子，优化器是否持有这些参数，以及每次更新后参数是否变化。只在这些关系成立之后，调整学习率和训练步数才有明确意义。

如果第一步 $A$ 的梯度为零而 $B$ 正常更新，这是零初始化设计的预期行为；如果很多步后两个因子都不动，则应检查是否同时置零、损失与分支脱离、梯度被清空得过早或学习率为零。若训练损失下降而生成结果没有变化，还需确认推理加载的是正确检查点，适配器被启用，生成模板相同，并检查任务指标是否对细微概率变化不敏感。

出现显存不足时，应先区分是权重加载、前向激活、反向峰值还是优化器更新阶段触发。减小微批次、缩短序列、使用重计算、量化基座或减少适配范围分别作用于不同部分。只把秩从较小值继续减半，可能仅节省很少状态，对长序列激活导致的瓶颈几乎没有帮助。

\subsection{怎样报告一个可信的结果}

质量指标应与任务一致：分类关注准确率及类别不平衡下的分项表现，结构化生成关注字段与约束正确性，代码任务关注在固定测试与执行环境下的通过率，开放回答需要明确评价规则与盲评设置。困惑度（perplexity）可定义为同一词元化与监督范围下平均负对数似然的指数，但跨分词器或不同掩码直接比较通常没有清楚含义。

资源报告应同时给出基座规模、可训练参数、权重与计算精度、序列长度、有效词元预算、硬件、峰值显存、吞吐和总时间。使用早停时，记录选择规则和实际停止位置。质量改善若伴随更长训练或更大模型，应把这种成本变化一并说明。适配器虽小，也可能记忆训练内容；发布前的数据权限与评估要求不会因为参数少而自动消失。

\section{综合练习与解答}
\label{sec:exercises}

\subsection{参数节省的成立条件}

\begin{example}
对一个 $m\times n$ 矩阵，只考虑原权重与两个因子，不计偏置。求 LoRA 参数数少于直接训练原矩阵的条件，并计算 $m=2048$、$n=4096$、$r=16$ 的情况。
\end{example}
\begin{solve}
条件为 $r(m+n)<mn$，即 $r<mn/(m+n)$。本例原矩阵有 $8388608$ 个参数，两个因子有
$16(2048+4096)=98304$ 个参数，比例为 $3/256\approx1.172\%$。参数节省取决于秩相对于两个维度是否足够小；若把秩提高到接近较小维度，因子参数量反而可能超过原矩阵。
\end{solve}

\subsection{零空间揭示了什么不能直接改变}

\begin{example}
固定一个 LoRA 层的 $A,B$，若输入 $x$ 满足 $Ax=0$，说明该层输出与基座输出的关系。进一步判断这是否意味着整个网络对对应原始样本一定没有变化。
\end{example}
\begin{solve}
该层的增量输出为 $sBAx=0$，因此对这一层收到的这个输入，输出等于基座线性层输出。但是更早适配器可能已经改变了到达该层的隐藏向量，更晚适配器也可能继续改变结果，所以不能推出整个网络输出不变。要讨论原始样本的端到端变化，必须沿完整计算图比较所有隐藏状态。
\end{solve}

\subsection{按词元累积梯度的权重}

\begin{example}
两个微批次分别有 $100$ 与 $300$ 个有效监督词元，其平均损失为 $L_1,L_2$。若希望优化全部 $400$ 个词元的平均损失，应如何组合？直接取 $(L_1+L_2)/2$ 会产生什么差异？
\end{example}
\begin{solve}
正确目标为 $(100L_1+300L_2)/400=L_1/4+3L_2/4$，其梯度使用相同系数组合。简单平均使第一微批次的总权重由四分之一增到二分之一，相当于其每个词元相对第二微批次获得三倍权重。该问题与是否使用 LoRA 无关，却会在比较不同批处理方案时改变训练目标。
\end{solve}

\subsection{合成秩与不能平均因子的反例}

\begin{example}
设 $B_1=(1,0)^{\mathsf T}$、$A_1=(1,0)$，
$B_2=(0,1)^{\mathsf T}$、$A_2=(0,1)$。求两项增量之和的秩，并比较直接相加与分别相加因子的结果。
\end{example}
\begin{solve}
$B_1A_1+B_2A_2=I_2$，秩为二，表明两个秩一更新之和不一定仍为秩一。另一方面，
$(B_1+B_2)(A_1+A_2)$ 是全一矩阵，秩为一，并出现原本不存在的非对角元素。若分别取平均，还会整体多出四分之一缩放，仍不等于所需的更新平均。
\end{solve}

\subsection{最优低秩增量的一个可解训练任务}

\begin{example}
设单层目标函数为
$\mathcal L(D)=\frac12\mathbb E\|(W_0+D)X-(W_0+M)X\|_2^2$，
其中 $\mathbb E[XX^{\mathsf T}]=I_3$、
$M=\operatorname{diag}(4,2,1)$，限制 $\operatorname{rank}(D)\leq1$。求最小损失，并说明为什么这个结果不保证任意因子 SGD 都能达到它。
\end{example}
\begin{solve}
基座相消后，目标为 $\frac12\|D-M\|_F^2$。由定理\ref{thm:truncated-svd}，最优可取
$D=\operatorname{diag}(4,0,0)$，最小损失为 $(2^2+1^2)/2=2.5$。
这是权重空间的全局最优值。因子参数化使目标对 $A,B$ 联合非凸；若两者均为零，前文已经构造出任务梯度为零而未达到最优的情形。学习率、初始化和停止条件也会影响有限步结果，所以表达集合中存在最优解与算法找到该解是两回事。
\end{solve}

\subsection{一个没有伪造结果的实验练习}

构造线性回归任务：选择 $n=m=16$，固定基座 $W_0$，生成标准正态输入，目标为
$Y=(W_0+M)X+\varepsilon$，其中噪声 $\varepsilon$ 均值为零、各维方差为给定的 $\sigma^2$，并独立于输入。分别令 $M$ 的奇异值迅速衰减与基本相等，比较秩为 $1,2,4,8,16$ 的适配器及全矩阵训练。训练集、验证集和测试集独立生成，并在各方法之间使用相同划分。

该实验能够区分三种误差来源：小秩无法表示目标造成的逼近误差，有限样本与噪声造成的估计误差，以及算法未充分优化造成的优化误差。在无噪声、无限数据的理想目标下，可用奇异值尾和给出理论最低误差；实际训练结果则需要运行后报告。应记录多个随机种子、参数量、训练曲线和测试误差，检查增加秩后是否真正靠近理论下界，而不是预先写下“高秩必然更好”的结论。

\section{进一步学习的路线与全文结论的层次}
\label{sec:closing}

完成本文后，第一层能力是准确检查维度、参数量、梯度和合并关系：看到一个适配实现，能够写出实际生效的 $W_0+sBA$，解释每个因子的形状，说明初始化时谁先获得梯度，并分辨哪些权重与状态真正被训练。第二层能力是理解优化与统计限制：低秩改变函数族和优化几何，较少参数既可能缓解过拟合，也可能限制表达；验证集表现必须与数据分布和训练预算一起解释。

更深入的理论研究可以学习矩阵流形优化、随机优化、谱正则化和局部二阶近似。若损失在某一点附近二阶连续可微，可以把参数扰动的影响近似写为线性项加由海森矩阵（Hessian matrix）确定的二次型；不同方向的曲率不同，说明相同权重误差可能带来不同损失变化。这一观点延续了式\eqref{eq:data-weighted-error}的数据加权思想，但一般深度网络 Hessian 可能非正定，且大步更新会离开局部近似范围，因此不能把二阶模型当作全局保证。

工程研究则应进一步关注长序列激活、量化算子、分布式分片、适配器路由和部署一致性。软件接口会演化，具体参数含义应以所用版本的官方文档为准，例如 PEFT 的 LoRA 配置文档\cite{peft-docs}。本文的稳定部分是数学关系与验证逻辑：先明确优化对象，再明确实现路径，最后用受控实验判断收益。低秩适配的价值在于给任务变化提供一种可管理的结构，而这种结构是否足够，始终由目标任务和证据决定。

\begin{thebibliography}{99}

\bibitem{hu2021}
Hu E. J., Shen Y., Wallis P., et al.
\emph{LoRA: Low-Rank Adaptation of Large Language Models}.
2021 年预印本，arXiv:2106.09685.
\url{https://arxiv.org/abs/2106.09685}.

\bibitem{vaswani2017}
Vaswani A., Shazeer N., Parmar N., et al.
\emph{Attention Is All You Need}.
2017，arXiv:1706.03762.
\url{https://arxiv.org/abs/1706.03762}.

\bibitem{dettmers2023}
Dettmers T., Pagnoni A., Holtzman A., Zettlemoyer L.
\emph{QLoRA: Efficient Finetuning of Quantized LLMs}.
2023，arXiv:2305.14314.
\url{https://arxiv.org/abs/2305.14314}.

\bibitem{kalajdzievski2023}
Kalajdzievski D.
\emph{A Rank Stabilization Scaling Factor for Fine-Tuning with LoRA}.
2023，arXiv:2312.03732.
\url{https://arxiv.org/abs/2312.03732}.

\bibitem{zhang2023}
Zhang Q., Chen M., Bukharin A., et al.
\emph{AdaLoRA: Adaptive Budget Allocation for Parameter-Efficient Fine-Tuning}.
2023，arXiv:2303.10512.
\url{https://arxiv.org/abs/2303.10512}.

\bibitem{liu2024}
Liu S.-Y., Wang C.-Y., Yin H., et al.
\emph{DoRA: Weight-Decomposed Low-Rank Adaptation}.
2024，arXiv:2402.09353.
\url{https://arxiv.org/abs/2402.09353}.

\bibitem{lora-code}
Microsoft.
\emph{LoRA: Official Implementation, loralib/layers.py}.
\url{https://github.com/microsoft/LoRA/blob/main/loralib/layers.py}.
访问日期：2026 年 10 月 3 日。

\bibitem{pytorch-autograd}
PyTorch.
\emph{Autograd Mechanics}.
\url{https://docs.pytorch.org/docs/main/notes/autograd.html}.
访问日期：2026 年 10 月 3 日。

\bibitem{pytorch-linear}
PyTorch.
\emph{torch.nn.functional.linear}.
\url{https://docs.pytorch.org/docs/main/generated/torch.nn.functional.linear.html}.
访问日期：2026 年 10 月 3 日。

\bibitem{peft-docs}
Hugging Face.
\emph{PEFT: LoRA Configuration and API Reference}.
\url{https://huggingface.co/docs/peft/main/package_reference/lora}.
访问日期：2026 年 10 月 3 日。

\end{thebibliography}
\end{document}
